\section{Introduction}

The United States built the interstate highway system in approximately 30 years, between the Federal-Aid Highway Act of 1956 and the completion of the final corridor segments in the mid-1980s. It has maintained that system poorly for the five decades since. The American Society of Civil Engineers' 2021 Infrastructure Report Card assigned US roads a C-minus; the Federal Highway Administration estimates that 45\% of the national highway network is in mediocre or poor condition \citep{FHWA_HPMS2022}. But the deterioration is not uniform. The corridors that carry the heaviest freight loads --- the T1 Primary Arteries identified in Track A of the ROUTE research module \citep{ROUTE_A1} --- have accumulated the largest reliability deficits, the highest climate exposure scores, and the most consequential structural gaps.

Interstate 2.0 is not a proposal to build a new network. It is a proposal to invest in the corridors that already carry most of the freight --- and to close the structural gaps, harden the vulnerable nodes, and integrate the emerging multimodal demand that the original 1956 design could not anticipate.

The ROUTE research module has produced five tracks of analysis supporting this synthesis:
\begin{itemize}
  \item \textbf{Track A} established a 15-dimension rubric for corridor classification, identifying 8 T1 Primary Arteries that together carry more than 50\% of national truck ton-miles \citep{ROUTE_A1,ROUTE_A2}.
  \item \textbf{Track B} identified 15 missing T1/T2 links, 11 compound resilience holes (corridors vulnerable on both climate and structural dimensions simultaneously), and found that 9 of 15 T1/T1 intersections have k=1 connectivity --- a single-point-of-failure condition \citep{ROUTE_B1,ROUTE_B3,ROUTE_B4}.
  \item \textbf{Track C} quantified the freight reliability deficit: \$8.2 billion per year in reliability costs on the NY--LA and HOU--CHI corridors alone, and demonstrated that I-69 completion would add 18\% capacity to the Gulf--Midwest max-flow \citep{ROUTE_C1,ROUTE_C2}.
  \item \textbf{Track D} mapped climate exposure and incident economics: \$6.2 billion per year in closure costs; the Donner Pass segment of I-80 as the single most consequential national disruption point; and a projection that Gulf Coast I-10 becomes the most critical climate adaptation investment by 2050 \citep{ROUTE_D1,ROUTE_D2}.
  \item \textbf{Track E} designed managed freight lanes for the seven highest-priority T1 corridors: \$121 billion capital, 2.3:1 benefit-cost ratio, with the potential to reduce transcontinental transit time by one full day \citep{ROUTE_E1}.
\end{itemize}

This paper synthesizes those tracks into a complete design framework and investment portfolio. The I2.0 program has three organizing principles:
\begin{enumerate}
  \item \textbf{Maximize throughput on existing T1 corridors before adding new ones.} Managed freight lanes and diamond interchange upgrades on the eight T1 arteries produce higher returns per dollar than new construction on any but the most deficient missing links.
  \item \textbf{Close the structural gaps identified by the rubric.} The gap map is not invented --- it is read from the dimension space of the scored corpus. Missing link construction targets corridors where the rubric score exceeds the T1 threshold but no T1-standard facility exists.
  \item \textbf{Harden compound exposure corridors against climate-driven failure.} The corridors that rank highest on both climate exposure (D1) and structural fragility (B1) receive targeted hardening investment before the 2050 risk window closes.
\end{enumerate}

The \$253 billion program cost deserves immediate context. The Infrastructure Investment and Jobs Act of 2021 allocated \$284 billion for surface transportation over five federal fiscal years --- approximately \$57 billion per year \citep{IIJA2021}. The I2.0 program, spread over 30 years, costs \$8.4 billion per year: comparable to, and likely substitutable within, the existing annual federal highway apportionment. This is not a call for new spending; it is a call for different priorities within the spending that already occurs.

\subsection*{Contributions}
\begin{enumerate}
  \item We synthesize five tracks of ROUTE analysis into the first complete, evidence-based investment portfolio for a targeted national interstate upgrade program.
  \item We demonstrate that the I2.0 investment portfolio has an aggregate NPV of \$298 billion at a 7\% discount rate (B/C ratio 2.2:1), with no single component below 1.2:1.
  \item We show that managed freight lanes on the highest-loaded T1 corridors produce a 50\% increase in commercial vehicle throughput and reduce transcontinental transit time from 4.5 to 3.5 days.
  \item We quantify the compound exposure hardening case: simultaneous closure of Donner Pass and I-35 (Oklahoma) saturates I-40 without hardening investment, producing network failure; with hardening, I-40 V/C reaches 0.91 --- stressed but stable.
  \item We demonstrate that a \$2 billion transit integration layer, added to the highway program at 0.8\% of total cost, unlocks 9 T1/T1 diamond hubs for intercity bus service, bringing an estimated 12.4 million transit-dependent Americans within 30 miles of an intercity connection they currently lack.
\end{enumerate}
